\chapter{Réalisation et implémentation}
\label{chap:realisation}

\callout{En bref}{Ce chapitre décrit la mise en œuvre concrète du système :
l'environnement technique, le profilage adaptatif, le calcul des indicateurs,
les visualisations (dont la carte géographique), les analyses avancées
(anomalies, segments), l'orchestration avec ses filets de sécurité, la
publication des livrables et les mécanismes de fiabilité.}

\section{Environnement et technologies}
Le système est développé en \textbf{Python~3}. Le tableau~\ref{tab:tech}
récapitule les principales technologies et leur rôle.

\begin{table}[ht]
\centering\small
\caption{Technologies utilisées.}
\label{tab:tech}
\begin{tabularx}{\linewidth}{@{}lY@{}}
\toprule
\textbf{Technologie} & \textbf{Rôle dans le projet} \\
\midrule
\textbf{pandas / NumPy} & Manipulation et analyse des données tabulaires. \\
\textbf{Plotly} & Génération des graphiques interactifs (HTML autonome). \\
\textbf{FastAPI / Uvicorn} & Interface web (téléversement, suivi, restitution). \\
\textbf{Jinja2} & Gabarits HTML de l'interface. \\
\textbf{\texttt{mcp}} & Implémentation du serveur Model Context Protocol. \\
\textbf{Groq / Ollama} & Fournisseurs de LLM (hébergé / local). \\
\textbf{kaleido} & Export statique des graphiques (figures du présent rapport). \\
\bottomrule
\end{tabularx}
\end{table}

L'organisation du code suit la séparation en couches : le paquet
\texttt{src/core/} contient le moteur déterministe, \texttt{src/llm/}
l'abstraction des fournisseurs de LLM, \texttt{src/mcp\_server/} le serveur
MCP et ses outils, \texttt{src/agents/} les sept agents, et \texttt{src/ui/}
l'interface web. Le pipeline (\texttt{src/pipeline.py}) relie l'ensemble.

\section{Interface utilisateur}
L'interface web, développée avec FastAPI et Jinja2, constitue le point
d'entrée. La figure~\ref{fig:accueil} présente la page d'accueil : zone de
téléversement, présentation des sept agents et historique des exécutions
passées.

\begin{figure}[ht]
\centering
\includegraphics[width=0.86\linewidth]{fig_accueil.png}
\caption{Page d'accueil : téléversement, présentation des agents et historique.}
\label{fig:accueil}
\end{figure}

Après téléversement, le traitement s'exécute en arrière-plan et l'utilisateur
suit sa \textbf{progression par étape}. La figure~\ref{fig:pipeline} montre la
vue d'exécution : statut de chaque agent, indicateurs de consommation (jetons,
coût estimé, latence) et résultat de la vérification des chiffres.

\begin{figure}[ht]
\centering
\includegraphics[width=0.74\linewidth]{fig_pipeline.png}
\caption{Suivi en temps réel de l'exécution du pipeline et synthèse des agents.}
\label{fig:pipeline}
\end{figure}

\section{Abstraction des fournisseurs de LLM}
Pour garantir la portabilité, l'accès aux modèles de langage passe par une
\textbf{abstraction} unique (\texttt{src/llm/client.py}) qui masque les
différences entre fournisseurs. Une fabrique sélectionne, selon la
configuration, le client \textbf{Groq} (hébergé, faible latence) ou
\textbf{Ollama} (local, sans clé d'API). Les deux exposent la même méthode
\texttt{chat()} acceptant des outils au format OpenAI et renvoyant un message
normalisé, enrichi de la consommation (jetons) et de la latence. Ce dernier
point alimente le suivi des coûts (section~\ref{sec:couts}).

\begin{lstlisting}[language=Python,caption={Sélection du fournisseur de LLM (\texttt{llm/client.py}).}]
def get_llm_client():
    provider = os.getenv("LLM_PROVIDER", "groq").lower()
    if provider == "groq":
        return GroqClient(api_key=os.getenv("GROQ_API_KEY"),
                          model=os.getenv("GROQ_MODEL"))
    if provider == "ollama":
        return OllamaClient(host=os.getenv("OLLAMA_HOST"),
                            model=os.getenv("OLLAMA_MODEL"))
    raise RuntimeError(f"Fournisseur inconnu : {provider!r}")
\end{lstlisting}

Cette conception applique le principe d'\textbf{inversion de dépendance} : les
agents dépendent d'une interface abstraite (\texttt{LLMClient}) et non d'un
fournisseur concret, ce qui permet d'en changer sans modifier le reste du
code.

\section{Ingestion, validation et profilage adaptatif}
L'agent d'\textbf{Ingestion} charge le fichier (l'extension détermine le
lecteur pandas approprié) puis valide sa structure (doublons, colonnes
entièrement vides, nombre de lignes). Le profilage, cœur de la généralité du
système, attribue à chaque colonne un rôle sémantique. L'extrait~\ref{lst:profiler}
illustre l'inférence du rôle pour les colonnes textuelles, où la détection
\textbf{géographique} est testée en priorité.

\begin{lstlisting}[language=Python,caption={Inférence de rôle pour une colonne textuelle (\texttt{profiler.py}).},label=lst:profiler]
# 4. Chaine de caracteres / objet
else:
    # La geographie est testee en premier et prime sur la cardinalite :
    # une colonne de noms de pays est geographique meme tres cardinale.
    if _looks_like_geo(s, geo_match_threshold):
        role = ColumnRole.GEO
    elif _looks_like_identifier(str(col), card, row_count, is_numeric=False):
        role = ColumnRole.IDENTIFIER
    elif card <= categorical_max_cardinality:
        role = ColumnRole.DIMENSION
    elif card > categorical_max_cardinality and card / max(row_count, 1) > 0.5:
        role = ColumnRole.TEXT
    else:
        role = ColumnRole.HIGH_CARD_DIM
\end{lstlisting}

La détection géographique compare les valeurs distinctes de la colonne à un
référentiel d'environ 250 noms de pays ; si une fraction suffisante
correspond (seuil paramétrable, 0,6 par défaut), la colonne devient une
dimension géographique exploitable par la carte.

\section{Calcul adaptatif des indicateurs (KPI)}
L'agent \textbf{KPI} récupère une liste généreuse de KPIs candidats déduits du
profil, en \textbf{cure} un sous-ensemble pertinent (jugement métier du LLM),
puis les calcule de façon déterministe. Chaque KPI est décrit par une
spécification JSON (identifiant, titre, type, mesure, agrégation, format), ce
qui permet un calcul générique. La figure~\ref{fig:kpis} montre les cartes
d'indicateurs telles qu'elles apparaissent dans le rapport.

\begin{figure}[ht]
\centering
\includegraphics[width=0.92\linewidth]{fig_rapport_kpis.png}
\caption{Indicateurs clés calculés automatiquement (extrait du rapport).}
\label{fig:kpis}
\end{figure}

\section{Visualisations et tableau de bord}
L'agent \textbf{Dashboard} sélectionne un ensemble \textbf{équilibré} de
graphiques (tendance, comparaison, distribution, composition, carte) et les
construit avec Plotly. Le moteur choisit automatiquement le type de graphique
selon les rôles des colonnes. La figure~\ref{fig:dashboard} présente le haut
du tableau de bord généré.

\begin{figure}[ht]
\centering
\includegraphics[width=0.9\linewidth]{fig_dashboard.png}
\caption{Tableau de bord interactif : navigation, indicateurs et début de la carte.}
\label{fig:dashboard}
\end{figure}

\subsection{La carte géographique (choroplèthe)}
Lorsqu'une colonne géographique est détectée, le moteur produit une
\textbf{carte choroplèthe} mondiale agrégeant une mesure par pays. La
figure~\ref{fig:carte} en donne un exemple : la distribution des buts à
domicile cumulés par pays.

\begin{figure}[ht]
\centering
\includegraphics[width=0.92\linewidth]{fig_carte.png}
\caption{Carte choroplèthe : buts à domicile cumulés par pays.}
\label{fig:carte}
\end{figure}

\subsection{Tendances temporelles}
Lorsqu'une colonne de type date est présente, le moteur génère des séries
temporelles, en agrégeant automatiquement par période adaptée au volume de
points. Les figures~\ref{fig:serieH} et~\ref{fig:serieA} illustrent
l'évolution des buts à domicile et à l'extérieur.

\begin{figure}[ht]
\centering
\begin{minipage}{0.49\linewidth}\centering
  \includegraphics[width=\linewidth]{fig_serie_home.png}
  \caption{Buts à domicile au fil du temps.}
  \label{fig:serieH}
\end{minipage}\hfill
\begin{minipage}{0.49\linewidth}\centering
  \includegraphics[width=\linewidth]{fig_serie_away.png}
  \caption{Buts à l'extérieur au fil du temps.}
  \label{fig:serieA}
\end{minipage}
\end{figure}

\subsection{Comparaisons et distributions}
Les graphiques en barres top-N classent les valeurs d'une dimension par mesure
(figures~\ref{fig:topT} et~\ref{fig:topE}), tandis que les
\textbf{histogrammes} révèlent la distribution des mesures
(figure~\ref{fig:distribution}), faisant apparaître l'asymétrie typique du
nombre de buts (forte concentration sur les faibles valeurs, longue traîne).

\begin{figure}[ht]
\centering
\includegraphics[width=0.62\linewidth]{fig_distribution.png}
\caption{Distribution des buts à domicile (histogramme).}
\label{fig:distribution}
\end{figure}

\begin{figure}[ht]
\centering
\begin{minipage}{0.49\linewidth}\centering
  \includegraphics[width=\linewidth]{fig_top_tournois.png}
  \caption{Top tournois par buts à domicile.}
  \label{fig:topT}
\end{minipage}\hfill
\begin{minipage}{0.49\linewidth}\centering
  \includegraphics[width=\linewidth]{fig_top_equipes.png}
  \caption{Top équipes par buts à domicile.}
  \label{fig:topE}
\end{minipage}
\end{figure}

\section{Analyses avancées : anomalies et segments}
L'agent KPI enrichit l'analyse de deux familles d'insights, implémentés dans
\texttt{src/core/insights.py}.

\subsection{Détection d'anomalies}
La détection repose sur un \textbf{score z modifié} fondé sur l'écart absolu
médian (MAD), robuste aux valeurs extrêmes que l'on cherche précisément à
isoler (extrait~\ref{lst:anomalies}).

\begin{lstlisting}[language=Python,caption={Score z modifié robuste (\texttt{insights.py}).},label=lst:anomalies]
def _modified_z(values):
    median = np.median(values)
    mad = np.median(np.abs(values - median))
    if mad == 0:
        return None  # pas de dispersion : repli sur l'ecart interquartile
    return 0.6745 * (values - median) / mad
\end{lstlisting}

Toute valeur dont le score dépasse un seuil (3,5 par défaut) est signalée. La
figure~\ref{fig:anomalies} montre la restitution dans le rapport.

\begin{figure}[ht]
\centering
\includegraphics[width=0.92\linewidth]{fig_rapport_anomalies.png}
\caption{Section « Anomalies » du rapport : valeurs aberrantes et exemple le plus extrême.}
\label{fig:anomalies}
\end{figure}

\subsection{Comparaison de segments}
Pour chaque dimension de faible cardinalité et chaque mesure, le moteur
compare les moyennes par groupe et retient les écarts les plus marqués
(figure~\ref{fig:segments}). Cette analyse met par exemple en évidence
l'avantage du terrain dans l'étude de cas.

\begin{figure}[ht]
\centering
\includegraphics[width=0.92\linewidth]{fig_rapport_segments.png}
\caption{Section « Comparaison de segments » du rapport.}
\label{fig:segments}
\end{figure}

\section{Orchestration et filets de sécurité}
Le pipeline (\texttt{src/pipeline.py}) exécute les sept agents dans un ordre
déterministe. Comme un agent peut, en pratique, omettre un appel d'outil, des
\textbf{filets de sécurité} déterministes vérifient après chaque étape que les
sorties attendues existent et, le cas échéant, invoquent eux-mêmes les outils
manquants sous les permissions du rôle concerné. L'extrait~\ref{lst:fallback}
illustre ce dispositif pour l'étape KPI.

\begin{lstlisting}[language=Python,caption={Filet de sécurité de l'étape KPI (\texttt{pipeline.py}).},label=lst:fallback]
elif stage == "kpi":
    if not store.get("kpis"):
        # ... calcul des KPIs si l'agent les a omis ...
        _safe("kpi", "compute_kpis", {"kpi_ids": top_ids})
    # Les insights sont independants des KPIs : on garantit chacun.
    if store.get("anomalies") is None:
        _safe("kpi", "detect_anomalies", {})
    if store.get("segments") is None:
        _safe("kpi", "compare_segments", {})
\end{lstlisting}

Ce mécanisme garantit qu'une exécution produit \textbf{toujours} des livrables
complets, même en cas de défaillance d'un agent — le moteur déterministe
devenant le filet de dernier recours.

\section{Publication des livrables}
L'agent de \textbf{Publication} assemble le tableau de bord interactif
(\texttt{dashboard.html}) et fige l'état complet de l'exécution
(\texttt{state.json}). L'agent \textbf{Reporter} produit ensuite le
\textbf{rapport décisionnel} (\texttt{report.html}), structuré en sections :
résumé, source des données, méthodologie, indicateurs, tableau de bord,
conclusions, anomalies, comparaison de segments. Le rapport est exportable en
\textbf{PDF} (impression navigateur) et en \textbf{HTML} autonome ; chaque
graphique est exportable en image. La figure~\ref{fig:entete} montre l'en-tête
du rapport avec sa barre d'actions et de navigation.

\begin{figure}[ht]
\centering

\includegraphics[width=0.92\linewidth]{fig_rapport_entete.png}
\caption{En-tête du rapport : navigation (Accueil, Dashboard) et export (PDF, HTML).}
\label{fig:entete}
\end{figure}

\section{Fiabilité : véracité des chiffres et suivi des coûts}

\subsection{Vérification des chiffres du rapport}
Pour contrer les hallucinations, chaque nombre cité dans les conclusions du
rapport est \textbf{vérifié} face à l'ensemble des nombres réellement calculés
(valeurs des KPIs, statistiques du profil, comptages, corrélations, valeurs
d'anomalies, moyennes de segments), avec une tolérance absorbant les arrondis.
Le rapport affiche alors un verdict : \emph{vérifié}, \emph{partiel} ou
\emph{non vérifié}. La figure~\ref{fig:findings} montre le badge de
vérification au-dessus des conclusions.

\begin{figure}[ht]
\centering

\includegraphics[width=0.92\linewidth]{fig_rapport_findings.png}
\caption{Conclusions du rapport avec le badge de vérification des chiffres.}
\label{fig:findings}
\end{figure}

\subsection{Suivi de la consommation}
\label{sec:couts}
Le client LLM mesure, pour chaque appel, le nombre de jetons (entrée/sortie)
et la latence ; le pipeline agrège ces mesures par agent et par exécution, et
estime un coût à partir de prix paramétrables. Ces indicateurs sont restitués
dans la vue d'exécution (figure~\ref{fig:pipeline}), assurant l'observabilité
du système.

\callout{Synthèse du chapitre}{Le système combine un moteur déterministe
généraliste (profilage, KPIs, graphiques, carte, insights) et des agents qui
curent et narrent. Les filets de sécurité garantissent des livrables complets,
la vérification des chiffres assure la véracité, et le suivi des coûts apporte
l'observabilité. Le chapitre suivant évalue l'ensemble.}
